% !TEX root = ../diary-book.tex

\section{令和7年01月02日 目標設定}

これまでは目標設定が曖昧でタスク分解できていなかったため、タスクの完了ができず、研究の進捗をはかることができていない。タスク分解してそれらの相関関係を作ることから始める。一方で、タスク分解ばかりに時間をかけることも意味がないため、随時アップデートできるような形にもっていく。

\href{https://xmind.ai/share/Cftre3Ed}{マインドマップ}を使って整理を進める。

2024年度の研究目的は基本設計を完了させることである。測位・時刻同期精度を一次宇宙線およびミュオンのエネルギーに対して見積もることが目標となる。

2025-01-02 09:25:43時点でのアウトラインは以下の通りである。

\begin{enumerate}
   \item 平面波近似での時刻・位置分解能の見積もり
   \item クライアントの時刻および位置の不定性による分解能の悪化についての検討（シャワー面同定頻度の影響）
   \item シャワー面同定頻度の見積もり
   \item ミュオン計測頻度の見積もり（深さ・物質厚依存性）
   \item ミュオンとシャワー面同時計測頻度の見積もり
\end{enumerate}